% ============================================================
% MÓDULO 7: DEEP Q-NETWORKS (DQN)
% ============================================================
\section{M7 · Deep Q-Networks}
\module{Módulo 7: Deep Q-Networks}{Sustituyendo la tabla por redes neuronales: estabilidad, replay buffer y target network}

\begin{frame}{El Salto a la Aproximación de Funciones}
  \begin{columns}[c,onlytextwidth]
    \column{0.50\textwidth}
      \textbf{OptiBot con Sensores Continuos:}
      \begin{itemize}
        \item Estado continuo $s \in \mathbb{R}^4$:
        \begin{enumerate}
          \item Distancia frontal LiDAR ($0.1$ a $15.0$ m)
          \item Porcentaje de batería ($0.0\%$ a $100.0\%$)
          \item Velocidad lineal ($0.0$ a $2.5$ m/s)
          \item Ángulo respecto a la estación de carga
        \end{enumerate}
        \item Hay infinitos estados posibles: \textbf{una tabla Q es inviable}.
        \item \textbf{Solución:} Una red neuronal paramétrica $Q(s, a; \theta) \approx Q^*(s, a)$.
      \end{itemize}
    \column{0.48\textwidth}
      \centering
      \begin{tikzpicture}[scale=0.85,transform shape]
        % Capa entrada
        \foreach \i in {1,2,3,4} {
          \node[circle,draw=Teal,fill=Teal!20,thick,minimum size=0.55cm] (in\i) at (0, 2.5 - 0.9*\i) {\tiny $s_\i$};
        }
        % Capa oculta
        \foreach \j in {1,2,3,4,5} {
          \node[circle,draw=Ink,fill=Mist!40,thick,minimum size=0.55cm] (h\j) at (2.0, 3.0 - 0.8*\j) {};
        }
        % Capa salida
        \foreach \k/\act in {1/Avanzar, 2/Girar, 3/Cargar} {
          \node[circle,draw=Coral,fill=Coral!20,thick,minimum size=0.55cm] (out\k) at (4.0, 2.1 - 0.9*\k) {};
          \node[right=0.1cm of out\k,font=\tiny] {$Q(s,\text{\act})$};
        }
        \foreach \i in {1,2,3,4} {
          \foreach \j in {1,2,3,4,5} {
            \draw[-{Latex[length=1.5mm]},SoftText!40] (in\i) -- (h\j);
          }
        }
        \foreach \j in {1,2,3,4,5} {
          \foreach \k in {1,2,3} {
            \draw[-{Latex[length=1.5mm]},Coral!40] (h\j) -- (out\k);
          }
        }
      \end{tikzpicture}
  \end{columns}
  \vspace{0.25cm}
  \keyidea{La red neuronal permite \textbf{generalizar}: estados similares obtienen valores Q similares aunque nunca antes hayan sido visitados.}
\end{frame}

\begin{frame}{La Tríada Mortal y los Dos Pilares de DQN}
  \begin{columns}[T,onlytextwidth]
    \column{0.46\textwidth}
      \begin{alertblock}{La Tríada Mortal (Sutton \& Barto)}
        Causa divergencia si se combinan:
        \begin{enumerate}
          \item \textbf{Aproximación de funciones}
          \item \textbf{Bootstrapping} (objetivo TD)
          \item \textbf{Entrenamiento Off-policy}
        \end{enumerate}
      \end{alertblock}
      \small Para evitar que la red diverja, DeepMind introdujo dos mecanismos estabilizadores fundamentales.
    \column{0.50\textwidth}
      \begin{block}{1. Experience Replay (Buffer $\mathcal{D}$)}
        \begin{itemize}
          \item Guarda tuplas pasadas $(s, a, r, s', d)$.
          \item Muestrea mini-batches aleatorios uniformes.
          \item Rompe la correlación temporal entre pasos consecutivos.
        \end{itemize}
      \end{block}
      \begin{block}{2. Target Network (Red Objetivo $\theta^-$)}
        \begin{itemize}
          \item Congela los pesos del objetivo durante $C$ pasos.
          \item Evita perseguir un ``blanco móvil''.
        \end{itemize}
      \end{block}
  \end{columns}
\end{frame}

\begin{frame}{Función de Pérdida y Double DQN}
  \begin{columns}[c,onlytextwidth]
    \column{0.52\textwidth}
      \textbf{Función de Pérdida de DQN:}
      \[
        \mathcal{L}(\theta) = \mathbb{E}_{(s,a,r,s',d) \sim \mathcal{D}} \left[ \Big( y_i - Q(s,a;\theta) \Big)^2 \right]
      \]
      con el objetivo de Bellman:
      \[
        y_i = r + \gamma (1 - d) \max_{a'} Q(s', a'; \mathbf{\theta^-})
      \]
    \column{0.45\textwidth}
      \begin{exampleblock}{Mejora: Double DQN}
        El operador $\max$ induce un \textbf{sesgo de sobreestimación} sistemático. Double DQN desacopla la decisión de la valoración:
        \[
          y_i^{\text{DDQN}} = r + \gamma Q\Big( s', \arg\max_{a'} Q(s', a'; \mathbf{\theta}); \mathbf{\theta^-} \Big)
        \]
        \begin{itemize}
          \item La red online $\theta$ elige la acción.
          \item La red target $\theta^-$ evalúa su valor.
        \end{itemize}
      \end{exampleblock}
  \end{columns}
\end{frame}

\begin{frame}[fragile]{Estructura de DQN en PyTorch}
  \begin{columns}[T,onlytextwidth]
    \column{0.64\textwidth}
\begin{lstlisting}[style=rlcode]
import torch
import torch.nn as nn

class QNetwork(nn.Module):
    def __init__(self, state_dim, action_dim):
        super().__init__()
        self.net = nn.Sequential(
            nn.Linear(state_dim, 64),
            nn.ReLU(),
            nn.Linear(64, 64),
            nn.ReLU(),
            nn.Linear(64, action_dim)
        )
    def forward(self, x):
        return self.net(x)

# Calculo de perdida con Replay Buffer
states, actions, rewards, next_states, dones = replay_buffer.sample(32)
with torch.no_grad():
    # Objetivo con Target Network congelada
    max_next_q = target_net(next_states).max(1)[0].unsqueeze(1)
    target = rewards + gamma * (1 - dones) * max_next_q

current_q = online_net(states).gather(1, actions)
loss = nn.MSELoss()(current_q, target)
\end{lstlisting}
    \column{0.33\textwidth}
      \begin{block}{Entrenamiento Estable}
        \begin{enumerate}
          \item Muestrear batch.
          \item Calcular target fijo.
          \item Backpropagation en $\theta$.
          \item Actualizar periódicamente $\theta^- \leftarrow \theta$.
        \end{enumerate}
      \end{block}
  \end{columns}
\end{frame}
